% Japanese and English cover pages, followed by Roman-numbered abstracts.
\begin{titlepage}
\thispagestyle{empty}
\centering
{\Large 2026年度 修士論文\par}
\vspace{2.6cm}
{\LARGE\bfseries システミックリスク分析に向けた\par}
\vspace{0.4cm}
{\LARGE\bfseries 動的銀行間ネットワーク形成\par}
\vfill
{\large 福岡工業大学大学院 工学研究科\par}
\vspace{0.35cm}
{\large システムマネジメント専攻\par}
\vspace{1.5cm}
{\large 学籍番号：JM24201\par}
\vspace{0.25cm}
{\large チョウ　コウウ\par}
{\large 張　浩宇\par}
\vspace{1.35cm}
{\large 指導教員：傅　靖\quad 准教授\par}
\vfill
{\large 提出日 2026年9月16日\par}
\end{titlepage}

\begin{titlepage}
\thispagestyle{empty}
\centering
{\Large 2026 Master's Thesis\par}
\vspace{2.2cm}
{\LARGE\bfseries Dynamic Interbank Network Formation\par}
\vspace{0.35cm}
{\LARGE\bfseries for Systemic Risk Analysis\par}
\vfill
{\large System Management Engineering\par}
{\large Master's Program\par}
{\large Graduate School of Engineering\par}
{\large Fukuoka Institute of Technology\par}
\vspace{1.5cm}
{\large ZHANG Haoyu\par}
{\large Student ID: JM24201\par}
\vspace{1.35cm}
{\large Under the Supervision of\par}
{\large Associate Professor Jing Fu\par}
\vfill
{\large September 16, 2026\par}
\end{titlepage}

% -----------------------------------------------------------------------------
% Japanese abstract page - arranged to follow the university sample layout.
% -----------------------------------------------------------------------------
\clearpage
\pagenumbering{roman}
\setcounter{page}{1}
\thispagestyle{plain}
\phantomsection
\addcontentsline{toc}{chapter}{要　旨}

\begingroup
\setlength{\parindent}{1em}
\setlength{\parskip}{0pt}

\begin{center}
{\fontsize{16pt}{24pt}\selectfont\bfseries
システミックリスク分析に向けた\par
動的銀行間ネットワーク形成\par}
\end{center}

\vspace{0.9cm}

\begin{flushright}
{\fontsize{10.5pt}{16pt}\selectfont
大学院工学研究科修士課程\par
システムマネジメント専攻\par
JM24201　張　浩宇\par}
\end{flushright}

\vspace{0.75cm}

\begin{center}
{\fontsize{14pt}{18pt}\selectfont\bfseries 要　旨\par}
\end{center}

\vspace{0.45cm}

{\fontsize{10.5pt}{17pt}\selectfont
銀行間市場は短期資金を配分する一方，取引によって形成された債権・債務関係を通じて損失を伝播させる可能性がある。しかし，既存研究では清算時のネットワークを所与とすることが多く，取引制度がエクスポージャーをどのように形成し，その後のリスクにどう影響するかは十分に明らかでない。本研究の目的は，集中型配分と分散型 RFQ の下で内生的に形成される銀行間ネットワークを比較し，資金配分，ネットワーク構造およびシステミックリスクの関係を検討することである。

本研究は，変化する銀行のバランスシート，内生的なネットワーク形成，期日付き契約，\mbox{Eisenberg--Noe} 清算，破綻処理および政策支援を一つの動的シミュレーションに統合する。集中型では一つの主体が市場全体の適格注文を観察し，共通の配分ルールを適用する。RFQ では各借り手が限定された局所的な貸し手集合に個別に照会し，提示された条件に基づいて取引を受け入れる。両方式には共通の会計，規制，契約および清算ルールを適用し，形成されたネットワークを通じて後の支払不足と資本不足がどのように生じるかを追跡する。

シミュレーションでは，実装された RFQ は集中型より多くの小口リンクを形成し，最大の二者間エクスポージャーと集中度を低下させた一方，資金充足度は低かった。基準設定では両方式の複合リスク経路は近く，RFQ の期間平均リスク差は統計的に有意ではないが，最終時点では RFQ のリスクと資本不足が低かった。契約条件と政策支援を変更すると差の大きさも変化した。したがって，一方の方式が常に優れているのではなく，ネットワーク分散化，資金充足度，資本不足および複合リスクの間にはトレードオフがあり，結論は実装された制度とパラメータに依存する。

\noindent\textbf{キーワード：}システミックリスク，銀行間ネットワーク，動的シミュレーション，RFQ，\mbox{Eisenberg--Noe} 清算
}

\vfill
\begin{flushright}
{\fontsize{10.5pt}{15pt}\selectfont 2026年9月16日}
\end{flushright}
\endgroup

% -----------------------------------------------------------------------------
% English abstract page - arranged to follow the university sample layout.
% -----------------------------------------------------------------------------
\clearpage
\thispagestyle{plain}
\phantomsection
\addcontentsline{toc}{chapter}{Abstract}

\begingroup
\setlength{\parindent}{1.5em}
\setlength{\parskip}{0pt}

\begin{center}
{\fontsize{18pt}{26pt}\selectfont\bfseries
Dynamic Interbank Network Formation\par
for Systemic Risk Analysis\par}
\end{center}

\vspace{0.9cm}

\begin{flushright}
{\fontsize{10.5pt}{16pt}\selectfont
JM24201 ZHANG Haoyu\par
System Management Engineering\par
Master's Program\par
Graduate School of Engineering\par}
\end{flushright}

\vspace{0.3cm}

\begin{center}
{\fontsize{14pt}{18pt}\selectfont\bfseries Abstract\par}
\end{center}

\vspace{0.15cm}
\begingroup
\fontsize{10.5pt}{17pt}\selectfont

Interbank markets allocate short-term funding, but the creditor--debtor links created by trading can also transmit distress. Existing network studies often take these links as given and therefore reveal less about how a trading mechanism creates the exposures that are later cleared. This thesis examines how centralized allocation and decentralized request-for-quote (RFQ) trading endogenously shape funding, network structure, and systemic risk.

The main contribution is an integrated dynamic framework connecting changing bank balance sheets, endogenous network formation, dated contracts, gross due-payment clearing under Eisenberg--Noe, absorbing default resolution, and policy support. In the centralized mechanism, one principal observes all eligible orders and applies a common allocation rule. In RFQ, borrowers search a limited local lender pool, request individual quotations, and accept bilateral trades without a market-wide coordinator. Both mechanisms operate under common accounting, regulatory, contract, and clearing rules, allowing the networks formed during trading to be followed into later payment and capital shortfalls.

The implemented RFQ mechanism forms more and smaller links and reduces both bilateral-exposure concentration and the largest individual exposure, but it also has lower funding satisfaction than the implemented centralized mechanism. Under the baseline specification, the composite-risk paths are close: the full-horizon mean difference is not statistically significant, although RFQ has lower risk and capital shortfall at the final observation. Changing contract and policy settings changes the magnitude of the comparison. The findings therefore identify a trade-off among network dispersion, funding completion, capital shortfall, and composite risk; they do not establish that either institutional design is universally superior.

\par
\noindent\textbf{Keywords:} systemic risk, interbank network, dynamic simulation,
decentralized request-for-quote, \mbox{Eisenberg--Noe} clearing

\par
\vspace{0.5em}

\begin{flushright}
\fontsize{10.5pt}{15pt}\selectfont
September 16, 2026
\end{flushright}

\endgroup
\endgroup

\clearpage
% Compile the complete thesis at least twice after final edits to refresh page numbers.
\begingroup
\hyphenpenalty=10000
\exhyphenpenalty=10000
\setlength{\emergencystretch}{3em}
\tableofcontents
\endgroup
\clearpage
